


\begin{ledgroupsized}[r]{120mm}
\footnotesize
\pstart
\noindent\textbf{Überlieferung:}
\pend
\end{ledgroupsized}
\begin{ledgroupsized}[r]{114mm}
\footnotesize
\pstart
\parindent -6mm
\makebox[6mm][l]{\textit{l}}%
Reinschrift
(Hand von J.\,D.~Brandshagen%
\indexN{\textls{Brandshagen} (Brondhaguen), Jobst Dietrich 1659\textendash nach 1716}) %
mit Verbesserungen von Leibniz (\textit{Lil}):
LH~XXXV~9,~20 Bl.~3\textendash4.
Ein Bogen~2\textsuperscript{o}.
Ränder ausgefranst.
\href{\Scan{LH_35_09_20_003wz}}%
{Wasserzeichen} in der Mitte von Bl.~3.
Eineinhalb Seiten auf Bl.~3~v\textsuperscript{o} und 4~r\textsuperscript{o}.
Bl.~3~r\textsuperscript{o} und das obere Drittel von Bl.~3~v\textsuperscript{o}
überliefern eine weitere Reinschrift durch Brandshagen:
Leibnizens Rezension zu \indexN{\textls{Vanni}, Giovanni Francesco 1638\textendash1709}\textsc{G.\,F.~Vanni},
\cite{02000}\textit{Exegeses physico-mathematicae}, Rom 1685, die als Textzeuge \textit{Lil} von N.\;\ref{41199} ediert wird.
Bl.~4~v\textsuperscript{o} bis auf den eigh. Ordnungsvermerk „Mathematica“ leer.
Der als Rezension konzipierte Text blieb unveröffentlicht.
\pend
\end{ledgroupsized}
\selectlanguage{latin}
\vspace{8mm}
\pstart
\centering
\href{\Scan{LH_35_09_20_003v}}%
{\lbrack3~v\textsuperscript{o}\rbrack}\ 
\hfill
\cite{02095}\textit{De momento quo gravia nituntur deorsum per datum planum};
\hfill
\phantom{\lbrack3~v\textsuperscript{o}\rbrack}
\linebreak 
\textit{quo subjectum planum urgent; quo planum inclinatum grave}\linebreak 
\textit{deorsum conatur; Dissertatio \indexN{\textls{Spoleti}, Francesco, gest. 1712}Francisci Spoleti.}\linebreak \textit{Altera ejusdem Dissertatio de secretione bilis in Hepate.} 
\edtext{Venet.~1686}{%
\lemma{Venet.}%
\Bfootnote{%
\textit{(1)}~1684
\textit{(2)}~1686~\textit{l}}}
in 4\textsuperscript{o}. 
\pend
\vspace{0.5\baselineskip}
\pstart\noindent
\edtext{\cite{02095}Libellus brevis at non inelegans 
\indexN{\textls{Christina}, Königin von Schweden 1632\textendash1654, gest. 1689}Christinae Reginae inscriptus est.}{%
\lemma{Libellus \lbrack...\rbrack\ est}%
\Cfootnote{\cite{02095}\textit{De momento quo gravia nituntur deorsum per datum planum} \lbrack...\rbrack\ \textit{de secretione bilis in Hepate}, Venedig 1686, S.\;3\,f.}}
\indexN{\textls{Spoleti}, Francesco, gest. 1712}Autor inter alia
\edtext{%
ostendit si globus plano inclinato\indexS{planum inclinatum} insistat\lbrack,\rbrack\indexS{globus plano inclinato insistens} momenta segmentorum seu 
\edtext{partium sphaerae\indexS{momenta partium spherae}}{%
\lemma{partium}%
\Bfootnote{%
\textit{(1)}~spherae~\textit{l}
\textit{(2)}~sphaerae~\textit{Lil}}}
citra ultraque punctum quo globus planum tangit\indexS{punctum quo globus planum tangit} 
consistentium differre 
momento quod sit ad momentum
\edtext{totalem sphaerae\indexS{momentum totale sphaerae}}{%
\lemma{totalem}%
\Bfootnote{%
\textit{(1)}~spherae~\textit{l}
\textit{(2)}~sphaerae~\textit{Lil}}}
ut perpendicularis ad hypotenusam inclinationis.\indexS{perpendicularis inclinationis}\indexS{hypothenusa inclinationis}
Cumque sphaera hic 
\edtext{unius segmenti}{%
\lemma{unius}%
\Bfootnote{%
\textit{(1)}~segmente\lbrack\textit{!}\rbrack~\textit{l}
\textit{(2)}~segmenti~\textit{Lil}}}
momento\indexS{momentum quo sphaera retrorsum ascendit} retrorsum ascendere, alterius 
\edtext{majoris vi\indexS{vis}}{%
\lemma{majoris}%
\Bfootnote{%
\textit{(1)}~qui~\textit{l}
\textit{(2)}~vi~\textit{Lil}}}
prorsum devolvi conetur; differentiam hanc momentorum\indexS{differentia momentorum} 
esse ipsam vim\indexS{vis qua sphaera devolvitur} statuit qua devolvitur 
\edlabel{LH_35_09_20_003-004_1a}%
\edtext{}{%
\xxref%
{LH_35_09_20_003-004_1a}{LH_35_09_20_003-004_1b}%
\lemma{sphaera.}%
\Bfootnote{%
\textit{(1)}~Quaeritur~\textit{l}
\textit{(2)}~Queritur~\textit{Lil}}}%
sphaera.%
}{\lemma{ostendit \lbrack...\rbrack\ sphaera}%
\Cfootnote{\cite{02095}a.\,a.\,O., pars~I, bes. S.\;17\,f.}}
\edtext{Queritur\edlabel{LH_35_09_20_003-004_1b}
autem aliquos demonstrationem suam infirmam proclamasse quasi Trigonometriae\indexS{Trigonometria}
et dimensioni circuli\indexS{dimensio circuli} inniteretur,%
}{\lemma{Queritur \lbrack...\rbrack\ inniteretur}%
\Cfootnote{\cite{02095}a.\,a.\,O., S.\;5.}}
sed testimonium perhibendum est veritati
fatendumque hanc dictorum momentorum differentiam\indexS{differentia momentorum} recte omnino sese habere,
nec levem Geometriae interioris\indexS{Geometria interior} peritiam 
\indexN{\textls{Spoleti}, Francesco, gest. 1712}autorem ostendisse.
\edtext{%
Detrahens autem momentum quo descendit sphaera,\indexS{momentum quo sphaera descendit} 
a momento ejus totali,\indexS{momentum totale}
reliquum esse ait momentum quo ipsa planum inclinatum premit.\indexS{momentum quo sphaera planum inclinatum premit}%
}{\lemma{Detrahens \lbrack...\rbrack\ premit}%
\Cfootnote{\cite{02095}a.\,a.\,O., S.\;19\,f.}}
\edtext{%
Denique si planum aliquod grave,\indexS{planum grave mobile} ponatur mobile circa axem, quo planum horizontis\indexS{planum horizontis} secat
planum inclinatum,\indexS{planum inclinatum} ait momentum quo planum mobile deorsum tendit\indexS{momentum quo planum mobile tendit deorsum} 
esse ad momentum totale ut sinum rectum anguli inclinationis\indexS{sinus rectus anguli inclinationis} ad sinum totum,\indexS{sinus totus}
idque pontibus versatilibus\indexS{pontes versatiles} applicat 
\edtext{qui arcium}{%
\lemma{qui}%
\Bfootnote{%
\textit{(1)}~artium~\textit{l}
\textit{(2)}~arcium~\textit{Lil}}}
custodiae apponi solent, quando ex 
\edtext{horizontali situ}{%
\lemma{horizontali}%
\Bfootnote{%
\textit{(1)}~sinu~\textit{l}
\textit{(2)}~situ~\textit{Lil}}}
circa cardinem versus verticalem debent attolli.%
}{\lemma{Denique \lbrack...\rbrack\ attolli}%
\Cfootnote{\cite{02095}a.\,a.\,O., S.\;21\,f.}}
\pend
\pstart
\href{\Scan{LH_35_09_20_004r}}%
{\lbrack4~r\textsuperscript{o}\rbrack}\ 
\edtext{Dissertatione posteriore statuit secretionem\indexS{secretio humorum corporis} humorum corporis\indexS{humor corporis} per solam cribrationem\indexS{cribratio} explicari non debere,
quasi diversae partes sanguinis\indexS{sanguis}\indexS{partes sanguinis} solo contactu jungerentur ut grana frumenti\indexS{grana frumenti} diversa,
sed cohaerere praeterea eo nexu, quo fluidorum villi\indexS{villi fluidorum} mutuo vinciuntur.
Hinc partes fluidorum\indexS{partes fluidorum} 
\edtext{minores minus}{%
\lemma{minores}%
\Bfootnote{%
\textit{(1)}~\textlangle\textendash\textrangle~\textit{l}
\textit{(2)}~minus~\textit{Lil}}}
facile se posse extricare.}{%
\lemma{Dissertatione \lbrack...\rbrack\ extricare}%
\Cfootnote{\cite{02095}a.\,a.\,O., S.\;24\,f.}}
\edtext{Considerat autem arterias\indexS{arteria} ut tubos conicos\lbrack,\rbrack\
semper enim magis in arctum per eas sanguis\indexS{sanguis} propellitur.}{%
\lemma{Considerat \lbrack...\rbrack\ propellitur}%
\Cfootnote{\cite{02095}a.\,a.\,O., suppositio VIII, S.\;27.}}
\edtext{Hinc putat 
\edtext{vincula fluidorum\indexS{vincula fluidorum}}{%
\lemma{}%
\Bfootnote{%
vincula \bars{illa \textit{gestr.~Lil}} fluidorum~\textit{l}}}
in illo ad angustias transitu magis dirumpi et alia ab aliis
\edtext{separari: et}{%
\lemma{separari}%
\Bfootnote{%
\textit{(1)}~.~Et~\textit{l}
\textit{(2)}~:~et~\textit{Lil}}}
ea quidem prius, quae minus cohaerint reliquis.}{%
\lemma{Hinc \lbrack...\rbrack\ reliquis}%
\Cfootnote{\cite{02095}a.\,a.\,O., suppositio V, S.\;27.}}
\edtext{Etsi autem Chymicorum\indexS{Chymici} fermentationem\indexS{fermentatio} pro Chymo\indexS{chymus} in ventriculo\indexS{ventriculus}
et Chylo\indexS{chylus} in Duodeno\indexS{duodenum} aliisque 
\edtext{agnoscat, negat}{%
\lemma{agnoscat,}%
\Bfootnote{%
\textit{(1)}~neget~\textit{l}
\textit{(2)}~negat~\textit{Lil}}}
tamen necessarium esse peculiare fermentum in parte secretioni\indexS{secretio} dicata.
Alioqui quaeri rursus 
\edtext{posset quonam}{%
\lemma{posset}%
\Bfootnote{%
\textit{(1)}~quinam~\textit{l}
\textit{(2)}~quonam~\textit{Lil}}}
priore fermento ipsum fermentum secretum fuerit.}{%
\lemma{Etsi \lbrack...\rbrack\ fuerit}%
\Cfootnote{\cite{02095}a.\,a.\,O., S.\;28.}}
\edtext{Quod speciatim ad secretionem bilis\indexS{bilis}\indexS{secretio bilis} 
attinet ostendit eam generari in hepate\indexS{hepar} ex sanguine venoso,\indexS{sanguis venosus}
Venamque portae\indexS{vena portae} ibi quodammodo arteriae\indexS{arteria} munere fungi, distribuendo sanguinem\indexS{sanguis} per innumera
capillamenta in glandulas secretorias\indexS{glandula secretoria} toto hepate\indexS{hepar} sparsas;
unde separatione facta bilis\indexS{bilis} quidem in ductum proprium\lbrack,\rbrack\
sanguis\indexS{sanguis} autem ea deposita
in venam cavam\indexS{vena cava} per innumera capillamenta feratur.}{%
\lemma{Quod \lbrack...\rbrack\ feratur}%
\Cfootnote{\cite{02095}a.\,a.\,O., S.\;30.}}
\edtext{Haec autem secretio\indexS{secretio} quo melius fiat in hepate,\indexS{hepar} ait sanguinem\indexS{sanguis} in lieni\indexS{lien}
semisolvi et disponi ad secretionem\indexS{secretio} in hepate\indexS{hepar} faciendam.
Animalia splenem\indexS{splen} execta\indexS{animal splenem exectum} salaciora esse et magis mingere,
quia supersit in sanguine\indexS{sanguis} aliquid execto liene\indexS{lien} quod ipsum reddat acriorem,
similiter in 
\edtext{mesaraicis\indexS{mesaraica} semisolvi}{%
\lemma{}%
\Bfootnote{%
mesaraicis \bars{\textlangle arteriis\textrangle~\textit{gestr.~Lil}} semisolvi~\textit{l}}}
bilem.\indexS{bilis} Porro semisolutum sanguinem\indexS{sanguis} cum bile\indexS{bilis} progredi in venam portae\indexS{vena portae}
cujus magna insuetaque latitudo, adeoque lentior in ea sanguinis motus,\indexS{motus sanguinis} 
spatium det partibus bilis semisolutis\indexS{partes bilis semisolutae} antea, sese invicem jungendi
ut fiant grandiores quod motus celer impediret,
quemadmodum ventus vapores impedire potest ne in pluviam coeant.}{%
\lemma{Haec \lbrack...\rbrack\ coeant}%
\Cfootnote{\cite{02095}a.\,a.\,O., S.\;31\,f.}}
\edtext{Hinc jam porro ex vena portae\indexS{vena portae} in angustiores canales intrusas cum sanguine\indexS{sanguis}
bilis partes in glandulis denique hepatis,\indexS{glandulae hepatis} prorsus secerni.}{%
\lemma{Hinc jam \lbrack...\rbrack\ secerni}%
\Cfootnote{\cite{02095}a.\,a.\,O., S.\;32.}}
\edtext{Caeterum negat bilem\indexS{bilis} continuo fluxu in intestina\indexS{intestinum} delabi, 
sed eo solum tempore quo refertus cibo ventriculus\indexS{ventriculus} vesicam felleam\indexS{vesica fellea} 
propinquam premat. Hinc in corporibus longo jejunio vexatis turgidum 
et multa repletum bile folliculum\indexS{folliculus} inveniri.}{%
\lemma{Caeterum \lbrack...\rbrack\ inveniri}%
\Cfootnote{\cite{02095}a.\,a.\,O., S.\;32.}}
\edlabel{LH_35_09_20_003-004_2a}%
Aliaque id genus habet \indexN{\textls{Spoleti}, Francesco, gest. 1712}autor, cujus eo magis laudandum est studium 
et ad majora excitandum quo pauciores in naturae animalis\indexS{natura animalis}
Mechanismis particulatim explicandis hactenus operae pretium fecerunt.%
\edlabel{LH_35_09_20_003-004_2b}%
\pend
